%!TEX root = ../dokumentation.tex

\chapter{Project Budget Summary}

Das Projektbudget gibt einen Überblick über die finanziellen Ressourcen,
die für die erfolgreiche Durchführung des Projekts erforderlich sind.
Gemäß dem Projektplanungsansatz nach Marchewka wird das Budget aus der
\ac{WBS}, den geschätzten Aktivitätsdauern sowie
den Kosten der den Aktivitäten zugeordneten Ressourcen abgeleitet.

Die Kostenschätzungen basieren auf den definierten Projektphasen sowie
den dafür vorgesehenen internen Projektressourcen. Die Aufwände für die
einzelnen Aktivitäten wurden mithilfe von \ac{PERT}-basierten
Dauerschätzungen ermittelt und anschließend unter Verwendung der
kalkulatorischen Tagessätze der jeweiligen Rollen in Kosten
überführt.

Das resultierende Projektbudget stellt somit die aggregierten Kosten
aller Projektaktivitäten dar, die zur Umsetzung des definierten
Projektumfangs erforderlich sind.

\begin{table}[htbp]
    \centering
    \caption{Projektbudget nach Projektphasen}
    \begin{tabular}{lr}
        \toprule
        \textbf{Projektphase}              & \textbf{Budget (€)} \\
        \midrule
        Projektdefinition                  & 16.000              \\
        Detailplanung \& Architekturdesign & 107.000             \\
        Systemkonfiguration \& Integration & 76.000              \\
        Test \& Qualitätssicherung         & 56.000              \\
        Pilotbetrieb                       & 98.000              \\
        Konzernweiter Rollout              & 55.000              \\
        Stabilisierung \& Projektabschluss & 36.000              \\
        \midrule
        \textbf{Gesamtbudget}              & \textbf{444.000}    \\
        \bottomrule
    \end{tabular}
\end{table}

Abbildung \ref{fig:budget_phase} visualisiert die Verteilung der
Projektkosten über die einzelnen Projektphasen.

\begin{figure}[htbp]
    \centering
    \includegraphics[width=0.9\textwidth]{costperphase.png}
    \caption{Projektkosten pro Phase}
    \label{fig:budget_phase}
\end{figure}

Das geschätzte Gesamtbudget des Projekts beträgt somit
\textbf{444.000 €}. In der Wirtschaftlichkeitsbetrachtung (\autoref{chap:projektbeschreibung}) wird mit Projektkosten von 500.000\,€ (inkl. Migration und Planungspuffer) gerechnet; die hier dargestellte Kostenbaseline entspricht der aus der \ac{WBS} abgeleiteten Planung. Die größten Budgetanteile entfallen auf die
Phasen \emph{Detailplanung und Architekturdesign} sowie
\emph{Pilotbetrieb}. In diesen Phasen fallen besonders hohe
Aufwände für Architekturentscheidungen, Systemkonfiguration,
Integration sowie die Validierung der Plattform im Rahmen eines
realitätsnahen Pilotbetriebs an.

Das dargestellte Budget stellt die \textbf{Kostenbaseline}
des Projekts dar. Während der Projektdurchführung werden die
tatsächlichen Projektkosten regelmäßig mit dieser Baseline
verglichen, um Abweichungen frühzeitig zu erkennen und
gegebenenfalls geeignete Steuerungsmaßnahmen einzuleiten.

Da die Umsetzung des Projekts ausschließlich mit internen
Ressourcen erfolgt, sind keine externen Implementierungs-
oder Beratungsleistungen im Projektbudget enthalten.